\color{unimain}{\subsection{Защита от непереключения фаз (ЗНФ)}}
\color{black}

\begin{enumerate}[label=\arabic{section}.\arabic{subsection}.\arabic{enumi}, labelsep=4pt, leftmargin=0pt, itemindent=57pt, itemsep=0pt, parsep=0pt] % если что можно крутить itemsep 0 или 5

\item
Защита от непереключения фаз (ЗНФ) предназначена для обнаружения расхождения полюсов выключателя с
пофазным приводом при его включении или отключении. 

\item
Ввод/вывод в работу ЗНФ осуществляется программным ключом \textbf{«ЗНФ Вn»}, либо с помощью сигнала «Вывод ЗНФ Вn».

\item
Факт расхождения полюсов может устанавливаться одним из следующих способов:
\begin{enumerate}[label=\arabic*)]
\item на основании внешнего сигнала от сборки блок-контактов выключателя;
\item функцией фиксации положения выключателя (ФПВ), которая пофазно контролирует состояние блок-контактов выключателя и в случае расхождения полюсов формирует сигнал \textit{«Внутр. пуск. ЗНФ Вn»}.
\end{enumerate}

\item
В зависимости от используемого способа соответствующий сигнал должен назначаться на
конфигурируемый вход \textit{«Пуск ЗНФ Вn»}.

\item
Действие на отключение выключателя осуществляется с выдержкой времени \textbf{«Тср ЗНФ Вn»}. Данная
уставка предназначена для отстройки от разновременности переключения блок-контактов выключателя.

\item
Предусмотрена возможность блокировки ЗНФ В1 от внешнего сигнала \textit{«Блок. ЗНФ В1»} на время, которое
задается уставкой \textbf{«Тдеблок ЗНФ В1»}.


\begin{figure}[H]
\centering
\includegraphics[width=1\textwidth,height=1\textheight,keepaspectratio]{\fbpath/05.02 ДЗО/761.3441 ЗНФ/img/ZNF.pdf}
\captionof{figure}{Функциональный блок ЗНФ В1}\label{pict_znf:ZNF}
\end{figure}


\small
\begin{longtable}{|>{\centering\arraybackslash}m{5.3cm}|>{\centering\arraybackslash}m{3.3cm}|>{\centering\arraybackslash}m{4.2cm}|>{\centering\arraybackslash}m{1.8cm}|>{\centering\arraybackslash}m{1cm}|}
\caption{Параметры для настройки ЗНФ\hfill\vspace{-0.5\baselineskip}}\label{znf:tbl1}\\ 
\hline
\rowcolor{gray!30}
Параметр & Условное обозначение на схеме & Значение/ Диапазон & Единица измерения & Шаг \\ 
\hline
\endfirsthead
\caption*{\hspace{3pt}\emph{Продолжение таблицы \ref{znf:tbl1}\hfill\vspace{-0.5\baselineskip}}} \\ % сделано по ГОСТ 2.105 п.6.8.7
\hline
\rowcolor{gray!30}
Параметр & Условное обозначение на схеме & Значение/ Диапазон & Единица измерения & Шаг \\ 
%\hline
\endhead

%===m>ЗНФ|-
%===m
\end{longtable}
\normalsize

\end{enumerate}